\section{Formal verification, reproducibility and transparency}
\label{sec:repro}

The release separates kernel checking, conventional proof, exact computation,
byte identity and measurement.  The formal project is the primary executable
evidence for the M0--M8 spine.  Exact integer and rational scripts remain
independent regression checks for the locked numerator, denominator, moment
boundary, arc controls and atlas invariants.  No floating-point tolerance
enters a theorem certificate.

\subsection{Lean build and trust report}

The formalisation pins both Lean and Mathlib to version 4.32.1.  Its root module
imports \path{M3A.AnalyticBridge}, which transitively imports every earlier
stage.  A clean \path{lake build} completed successfully after M8.  The
principal axiom reports contain only \path{propext},
\path{Classical.choice} and \path{Quot.sound}; there is no
\path{sorry}, \path{admit}, user-declared axiom or unsafe theorem dependency.
The theorem-level map is reproduced in \cref{app:reprochecks}.

This trust statement follows the standard Lean architecture
\parencite{deMouraUllrich2021}: tactics and automation may construct proof
terms, but acceptance of those terms is checked by the kernel.  Mathlib
provides the matrix, power-series, topology and convexity infrastructure used
in the development \parencite{MathlibCommunity2020}.  Neither citation is used
as evidence for the M3A theorems themselves; their evidence is the released
source and checked proof terms.

\subsection{Independent exact routes}

\begin{enumerate}[label=\arabic*.]
\item \textbf{Lean route.}  Finite matrix definitions, cofactor expansion,
      word recurrence, support classification, convex hull and analytic
      summation form one checked dependency chain.
\item \textbf{Cofactor route.}  A direct Leibniz determinant gives (N) and the
      52-term denominator.
\item \textbf{Moment route.}  Polynomial matrix powers give $M_k$, endpoint
      support and the signed minimal layer.
\item \textbf{Pulled-amplitude route.}  Gaussian-rational series compute the
      ARC1 and GF1 tie cancellations and next survivors.
\item \textbf{Combinatorial route.}  Released JSON atlases are checked for all
      15 MV1 cells and all 16 GF1 cones, including radicality flags and Hilbert
      data.
\end{enumerate}
The cofactor and moment routes meet at the Laurent expansion of the resolvent;
agreement is a cross-check, not circular evidence.  Although finite
Hamiltonian entries generate highly structured series, no general D-finite
closure claim is needed here \parencite{Stanley1980}.

\subsection{Executable discrepancy fixture}

The fixture changes the symmetric edge pair
$H_{0,2}=H_{2,0}=-1$ to $+1$.  It then reruns the same cofactor and moment
code.  Both routes detect disagreement with the locked numerator and moment
layer.  The fixture result is computed at run time and written to
\path{DATA/discrepancy_detection.json}; it is not a stored truth value.

\subsection{Source recovery}

The workstation scan hashed 3,636 scientific-source files and found 932 unique
byte streams, including 631 duplicate-content groups.  It searched 780 unique
text streams and extracted 20 unique PDFs without a read failure.  Those counts
are \evidence{MEASURED} with the stated sample sizes.  Vendor TeX
distributions, dependency caches, rendered scratch files and the output tree
were excluded and recorded.  Byte-identical copies were searched once and
linked by SHA-256; filename suffixes such as ``copy'' or parenthetical numbers
were not treated as authority.

\paperone{} has the verified Zenodo version DOI
\href{https://doi.org/10.5281/zenodo.20556571}{10.5281/zenodo.20556571} and is
cited accordingly \parencite{MedfordPaperI}.  Its source is available in the
\href{https://github.com/SLresearch74/Signed-Involution-Sector-Exclusion-for-EZT-in-Finite-Magnetic-Graph-Hamiltonians---v1.0.0}{public GitHub repository}.
The Zenodo record is version \path{v1.0.0}; as of 10 August 2026 the GitHub
repository has no separate Release or tag.  Paper II reproducibility materials
are located under
\path{subsequent-work/paper-2-opening-theory/} in that repository.  The earlier
R1.2 Paper II manuscript remains an author-supplied unpublished preprint with
no invented identifier.

\subsection{How to reproduce}

From the formalisation root, run
\begin{verbatim}
lake build
lake env lean M3AFormalisation.lean
\end{verbatim}
From \path{subsequent-work/paper-2-opening-theory}, run
\begin{verbatim}
python CODE/verify_paper.py
python CODE/verify_paper.py --fixture
cd paper/source
latexmk -pdf -interaction=nonstopmode -halt-on-error main.tex
\end{verbatim}
The paper build uses Biber for Harvard author--date citations.  Exact
environment versions, milestone verification reports, expected JSON and the
release manifest accompany the source.  The full denominator, atlases and
claim maps are machine-readable rather than transcribed by hand.

\subsection{Known reproducibility boundary}

The current Lean package is complete through the fixed-fibre analytic theorem,
but not through the paper's later pullback, arc, atlas, Gr\"obner/Rees or
Noetherian sections.  Their exact verifiers check formulas, examples and finite
atlases; they do not turn those conventional proofs into kernel-checked ones.
Separately, the fixed-weight theorem proves finite generation rather than an
effective uniform bound.  Its verifier checks hypotheses and examples but
cannot turn a non-effective Noetherian argument into a global finite atlas.
Both limitations are reported rather than hidden.
